% !TEX root = ../../../main.tex

\subsection{特征数据集构建与输入组织}

句子视频识别模型的训练数据来自 WLASL 小词表中的词级视频。
每个原始视频通常对应一个独立手语词汇，因此训练阶段的目标是将每个词级视频转换为一个固定形状的时序特征样本。
与实时单词识别中的自采 raw dataset 不同，句子视频识别实验直接以离线视频文件作为输入。
系统需要先从视频中提取人体姿态和双手关键点，再进一步组织为模型可以训练的特征数组。

\WhuAutoFigure{0.82\textwidth}{0.42\textheight}{figures/chapter06/sentence-feature-build-flow.png}{句子视频识别特征数据集构建流程图}{fig:chapter06-sentence-feature-build-flow}

如图~\ref{fig:chapter06-sentence-feature-build-flow} 所示，句子视频识别特征数据集的构建过程可以划分为视频读取、关键点检测、动作窗口定位、帧重采样、特征构造和数据集输出几个步骤。
系统首先根据样本清单读取每个视频的文件路径和 Gloss 标签。
其中，视频路径用于定位原始 WLASL 词级视频文件，Gloss 标签用于确定该视频对应的词汇类别。
随后，系统使用 OpenCV 读取视频帧，并将每一帧送入 MediaPipe Holistic 模型进行处理。
Holistic 检测结果中包含人体姿态关键点、左手关键点和右手关键点，为后续 plus 特征构造提供原始点位来源。

在视频帧处理阶段，系统不会直接从整段视频中等间隔抽取 20 帧。
这是因为词级视频中可能包含动作开始前的静止帧、动作结束后的停顿帧以及手部尚未进入主要动作区域的过渡帧。
如果直接对整段视频进行均匀采样，采样结果可能包含较多无效帧，从而削弱模型对核心动作片段的学习效果。
因此，系统首先根据手部检测结果确定有效动作窗口，再在动作窗口内部进行重采样。

为了定位有效动作窗口，系统会为视频中的每一帧生成 hand\_flags 标记。
当某一帧检测到左手或右手关键点时，该帧的 hand\_flags 记为 True。
当左右手均未被检测到时，该帧的 hand\_flags 记为 False。
系统统计 hand\_flags=True 的帧数，用于判断该视频中是否存在足够可靠的手部动作区域。
如果有效手部帧数量达到设定阈值，系统以第一个 hand\_flags=True 的帧作为动作区域起点，以最后一个 hand\_flags=True 的帧作为动作区域终点。
随后，系统会在动作区域前后各扩展一定数量的 padding 帧，以保留动作进入和动作收尾阶段的信息。

当有效手部帧数量不足时，系统不会直接丢弃该视频。
为了保证数据集构建过程具有一定的容错能力，系统会退化为使用视频中部窗口作为采样区域。
这种回退策略可以避免由于个别帧关键点检测失败而导致样本完全无法使用。
同时，样本来源、窗口位置和采样帧信息会被记录到 sample\_index.csv 中，便于后续进行数据质量检查和问题样本追踪。

动作窗口确定后，系统从窗口范围内重采样 20 个帧下标。
重采样采用线性采样方式，从动作窗口起点到终点均匀生成 20 个采样位置，并将采样位置转换为实际帧下标。
通过这种方式，不同长度的视频片段都可以被转换为统一长度的时序输入。
对于持续时间较短的视频，重采样可以保留其主要动作变化。
对于持续时间较长的视频，重采样可以压缩冗余帧，使模型重点学习动作的整体变化趋势。

对于每一个采样帧，系统调用 plus 特征构造方法生成单帧 232 维特征。
该特征由 base\_166、static\_plus\_28 和 dynamic\_delta\_38 三部分组成。
其中，base\_166 用于描述左右手手型结构和上肢基础姿态。
static\_plus\_28 用于补充单帧内的人体相对位置、手部中心、手部范围、手部存在状态和距离信息。
dynamic\_delta\_38 用于描述相邻帧之间姿态点和手部代表点的变化趋势。
因此，一个词级视频样本最终会被组织为 $20\times232$ 的时序特征矩阵。

完成单个视频的特征构造后，系统将该样本的特征矩阵加入输入数组 $X$。
同时，系统根据该视频对应的 Gloss 标签生成类别编号，并将其加入标签数组 $y$。
在处理全部视频样本后，系统将所有样本统一保存为特征数据集文件。
其中，$X$ 保存模型训练所需的时序特征，$y$ 保存每个样本对应的类别编号。
标签映射文件用于记录类别编号与 Gloss 标签之间的对应关系，保证模型推理结果能够反查到具体词汇。

\begin{longtable}{L{0.24\textwidth} L{0.32\textwidth} L{0.34\textwidth}}
  \caption{句子视频识别特征数据集文件说明表}
  \label{tab:chapter06-sentence-feature-files}\\
  \toprule
  文件名 & 内容 & 作用 \\
  \midrule
  X.npy & 样本数 $\times$ 20 $\times$ 232 特征数组 & 模型输入 \\
  y.npy & 类别编号数组 & 模型标签 \\
  labels.json & Gloss 标签映射 & 推理结果反查 Gloss \\
  sample\_index.csv & 样本来源、动作窗口、采样帧等信息 & 数据追踪与质量检查 \\
  config.json & 特征构建配置 & 记录构建参数 \\
  \bottomrule
\end{longtable}

如表~\ref{tab:chapter06-sentence-feature-files} 所示，特征数据集由模型输入、标签、标签映射、样本索引和配置文件共同组成。
其中，X.npy 和 y.npy 是模型训练阶段直接读取的核心文件。
labels.json 保存 Gloss 标签与类别编号的映射关系，用于训练后解释模型输出结果。
sample\_index.csv 记录每个样本对应的原始视频、动作窗口范围和实际采样帧下标。
该文件可以帮助开发者检查某个样本是否存在动作窗口定位异常、有效帧不足或采样位置不合理等问题。
config.json 则保存本次特征构建使用的参数，例如目标帧数、特征维度、动作窗口 padding 和最小有效手部帧阈值等。

通过上述处理，系统将来源不同、长度不同的视频样本统一转换为固定形状的时序特征数据。
这种输入组织方式既便于后续双向 GRU 时序分类模型进行批量训练，也为句子视频识别中的滑动窗口预测和候选词输出提供了统一的数据格式基础。
同时，sample\_index.csv 和 config.json 的保存也提高了实验的可复现性，使后续模型效果分析和特征构建参数调整更加方便。